\section{Related Work}
\label{sec:related_work}

\textbf{Static Model Merging.}
Model merging creates unified multi-task networks from fine-tuned experts without retraining. Static approaches like Model Soups \cite{wortsman2022model} and Git Re-Basin \cite{ainsworth2022git} average fine-tuned weights to improve generalization. However, on distinct, non-overlapping tasks, simple averaging suffers from task interference. To alleviate this, Task Arithmetic \cite{ilharco2022editing} adds scaled task vectors. TIES-Merging \cite{yadav2023ties} addresses sign conflicts and parameter redundancies by filtering low-magnitude updates and resolving sign agreement. DARE \cite{yu2024dare} uses stochastic dropping and rescaling to sparsify task vectors. While static methods are computationally free during inference, they cannot adapt to dynamic input-level task shifts.

\textbf{Dynamic Model Merging.}
To overcome the capacity constraints of static merging, input-adaptive dynamic model merging has been proposed. AdaMerging \cite{yang2023adamerging} optimizes merging coefficients on a tiny unlabeled dataset at test time, but its coefficients remain static during subsequent deployment. 
To achieve input-dependent merging, methods like Linear Routers and QWS-Merge compute dynamic coefficients using a routing network applied to the input batch. 
However, as analyzed in Section~\ref{sec:intro}, these methods reconstruct full-scale dense weights *per batch* of inputs by averaging routing coefficients across the batch dimension. This introduces a fatal batch-size dependency, results in catastrophic \emph{heterogeneity collapse} when heterogeneous tasks are mixed in a single batch, and incurs prohibitive execution and memory costs to rebuild dense weights during the forward pass. Our proposed SLD-Merge completely departs from this paradigm by shifting dynamic routing from weight reconstruction to sample-wise activation-space routing.

\textbf{Low-Rank Decomposition and Mixture of Experts.}
Parameter-Efficient Fine-Tuning (PEFT) methods, such as LoRA \cite{dettmers2024qlora}, adapt large pre-trained models by training low-rank matrices. Our approach is conceptually related but operates in a post-hoc, training-free manner: instead of training low-rank adapters from scratch, we perform offline Singular Value Decomposition (SVD) on already fine-tuned, fully dense task vectors of specialized experts, keeping only the top-$r$ singular vectors. This is also related to Mixture-of-Experts (MoE) \cite{shazeer2017outrageously, fedus2022switch}, which routes tokens through specialized feed-forward networks. However, standard MoE models are trained from scratch under massive compute budgets, whereas SLD-Merge is a lightweight, training-free merging framework designed to consolidate existing dense expert models.
By combining SVD with Top-1 hard gating, SLD-Merge achieves the high multi-task performance of dynamic routing while preserving the extremely low storage and compute footprint required for real-world edge deployment.
